% Generated draft: requires mathematical and source review.
\documentclass[11pt]{article}
\usepackage[margin=25mm]{geometry}
\usepackage{fontspec}
\setmainfont{Latin Modern Roman}
\usepackage{amsmath,amssymb,mathtools,unicode-math}
\setmathfont{Latin Modern Math}
\usepackage{enumitem,needspace,xurl,hyperref}
\setlength{\parindent}{0pt}
\setlength{\parskip}{5pt}
\setlength{\emergencystretch}{3em}
\newcommand{\N}{\mathbb{N}}
\newcommand{\Z}{\mathbb{Z}}
\newcommand{\Q}{\mathbb{Q}}
\newcommand{\R}{\mathbb{R}}
\newcommand{\C}{\mathbb{C}}
\newcommand{\F}{\mathbb{F}}
\newcommand{\A}{\mathbb{A}}
\newcommand{\PP}{\mathbb{P}}
\newcommand{\E}{\mathbb{E}}
\newcommand{\eps}{\varepsilon}
\newcommand{\dd}{\,\mathrm{d}}
\begin{document}
\section*{UnsolvedMath ID: 30003925}
\textbf{Source problem code:} catalogue entry
\subsection{Definitions and mathematical statement}
A domain is a nonempty connected open set. For a bounded domain \(\Omega\subset\C^2\), put \(K=\overline{\Omega}\). Pseudoconvexity means that \(\Omega\) admits a continuous plurisubharmonic exhaustion function. For \(r\in\{2,\infty\}\), boundary regularity of class \(C^r\) means that near each boundary point there is a real-valued \(C^r\) defining function \(\rho\) with \(\Omega=\{\rho<0\}\) locally and \(\dd\rho\ne0\) on the boundary.

Using the usual interpretation that the neighborhood basis concerns the closure, \(K\) has a Stein neighborhood basis if, for every open set \(U\subset\C^2\) containing \(K\), there is a Stein open set \(V\subset\C^2\) such that
\[
K\subset V\subset U.
\]
The supplied record does not explicitly specify that the basis concerns \(\overline{\Omega}\); this interpretation requires confirmation against the cited report.

\textbf{Existence question.} For each of the regularity classes \(C^2\) and \(C^\infty\), does there exist a bounded pseudoconvex domain \(\Omega\subset\C^2\) whose boundary has that regularity, whose closure has a Stein neighborhood basis, and which is not \(1\)-\(H\)-convex in the sense intended by the cited source?

\textbf{Unresolved qualification.} The supplied record gives no definition of \(1\)-\(H\)-convexity, including its underlying object and any quantitative neighborhood conditions. That term is retained as a source-dependent predicate; its negation cannot be expanded into a precise mathematical condition from the supplied evidence alone. Thus the existence question above remains conditional on recovering that definition. The source writes \(C^2\) ``or'' \(C^\infty\); the two regularity variants are displayed separately here. An example with \(C^\infty\) boundary would also satisfy the \(C^2\) requirement, while the converse does not follow from the stated regularity.

The supplied dated assessment reports a \(C^{1,1}\) counterexample and labels the \(C^2\) and \(C^\infty\) existence question unresolved. These are supplied literature claims, not independently verified mathematical conclusions.
\subsection{Short English statement}
The question asks whether a bounded pseudoconvex domain in \(\C^2\) can have a Stein neighborhood basis and fail the source's \(1\)-\(H\)-convexity condition when its boundary is \(C^2\), or even \(C^\infty\). The missing definition of \(1\)-\(H\)-convexity prevents a fully explicit formulation. The supplied curation reports only a \(C^{1,1}\) counterexample and labels the smoother cases open.
\subsection{Sources}
\begin{itemize}
\item Local UnsolvedMath registry, supplied complete record: permanent ID \(30003925\), source problem code \texttt{catalogue entry}, in the Oberwolfach Reports Open Problems collection. No individual proposer is supplied; the proposed year is recorded as 2018.
\item Supplied source citation: \emph{Geometric Methods of Complex Analysis (2019)}, Oberwolfach Reports, pp.\ 2253--2302, DOI \texttt{10.4171/owr/2018/37}. Supplied source and literature link, labeled ``Oberwolfach Report 37/2018'': \url{https://doi.org/10.4171/owr/2018/37}. The supplied background describes Oberwolfach Reports as documenting workshops and their open-problem discussions.
\item Supplied literature reference: \emph{Strong Stein neighborhood bases}, \url{https://arxiv.org/abs/0705.0507}. The record attributes a \(C^{1,1}\) counterexample to cited work without identifying a theorem, page, or precise construction.
\item Literature assessment dated exactly \texttt{2026-08-21}: the supplied curation labels the problem open and reports that the \(C^2\) or \(C^\infty\) existence question remains unresolved despite a \(C^{1,1}\) counterexample. This is supplied dated curation derived from the OWR record, not an independently verified literature assessment. Neither linked source has been independently consulted for this draft.
\end{itemize}
\end{document}
